\section{Diskussion}

Dieses Kapitel ordnet die zuvor berichteten quantitativen und qualitativen Befunde inhaltlich ein. Zunächst erfolgt eine detaillierte Beantwortung der beiden zentralen Forschungsfragen. Daran anschließend wird eine methodische Reflexion durchgeführt, welche insbesondere die Diskrepanz zwischen den verwendeten Evaluationsmetriken und dem eigentlichen didaktischen Ziel beleuchtet. Abschließend werden die Auswirkungen der Datenqualität sowie die methodischen Limitationen der vorliegenden Studie kritisch diskutiert.

\subsection{Beantwortung von RQ1 und RQ2}

Die Auswertung der ersten Forschungsfrage zeigt unmissverständlich auf, dass ein Sprachmodell ohne spezifische Steuerung nicht als sokratischer Tutor agieren kann. Die ungepromptete Referenz liefert auf fachliche Fragen sofort vollständige Erklärungen und direkte Lösungswege. Erst durch die Implementierung eines dedizierten Prompts ändert sich dieses Verhalten grundlegend. Dabei erweist sich die ausführliche Variante in Form des Systemprompts als die robusteste und effektivste Lösung. Diese Variante hält die finale Antwort am zuverlässigsten zurück und zwingt die Lernenden durch beständig weiterführende Fragen zum eigenen Nachdenken.

Die Evaluation offenbart hierbei jedoch auch einen immanenten didaktischen Zielkonflikt der sokratischen Methode. Wenn der Tutor einer hartnäckigen Forderung nach der direkten Lösung konsequent widersteht, wird dies von den Auswertungsmetriken fälschlicherweise als mangelnde Anpassung an das Niveau interpretiert. Der scheinbare Leistungsabfall der strengsten Variante in diesem spezifischen Kriterium ist somit kein Beweis für ein schlechtes System, sondern belegt paradoxerweise gerade die exakte Einhaltung der didaktischen Vorgaben.

Bei der zweiten Forschungsfrage zum Vergleich der unterschiedlichen RAG Konfigurationen zeichnet sich ein weitaus weniger klares Bild ab. Zwar deuten die erhobenen Daten auf graduelle Qualitätsunterschiede zwischen den getesteten Retrieval Ansätzen hin, doch lassen sich diese Effekte kaum isoliert betrachten. Die im späteren Verlauf detailliert beschriebene geringe Qualität der zugrunde liegenden Datenbasis überlagert die potenziellen Unterschiede der Konfigurationen nahezu vollständig. Eine abschließende und allgemeingültige Beantwortung dieser zweiten Frage ist daher auf Basis der vorliegenden Ergebnisse nur sehr eingeschränkt möglich.

\subsection{Metrik Mismatch (Goal Accuracy und Completeness vs sokratisches Ziel) und methodische Reflexion}

Die automatisierte Auswertung tutorieller Dialoge durch generische Metriken bringt eine erhebliche methodische Herausforderung mit sich. Der vorliegende Datensatz deckt einen massiven Metrik Mismatch auf. Klassische Kriterien der Textgenerierung wie Goal Accuracy oder Completeness sind primär darauf ausgelegt, eine möglichst schnelle, präzise und vollumfängliche Beantwortung einer Nutzeranfrage positiv zu bewerten. Das sokratische Lehrziel verlangt hingegen exakt das Gegenteil. Ein guter Tutor gibt die Lösung eben nicht direkt vor, sondern leitet die studierende Person durch kleine kognitive Hürden und strukturierte Rückfragen zur selbstständigen Erarbeitung des Wissens.

Diese bewusste Lösungszurückhaltung wird von herkömmlichen Metriken paradoxerweise als unvollständige oder ausweichende Antwort bestraft. Besonders deutlich wird dies in Dialogen, in denen die simulierten Lernenden explizit nach dem Endergebnis verlangen. Bleibt das System hierbei in seiner sokratischen Rolle, sinken die Scores für die allgemeine Hilfsbereitschaft und die Niveauanpassung drastisch ab. Folglich messen die verwendeten Standardmetriken teilweise nicht die tatsächliche didaktische Qualität, sondern bevorzugen vielmehr einen dozierenden Erklärstil. Diese methodische Diskrepanz muss bei zukünftigen Evaluationen durch speziell angepasste sokratische Bewertungsmaßstäbe zwingend behoben werden, um den wahren Wert eines Tutorsystems abbilden zu können.

\subsection{Datenqualität (leere Gespräche und 30\,\% Retrieval Abdeckung) und deren Einfluss}

Ein zentraler Faktor bei der Interpretation der Gesamtergebnisse ist die stark limitierte Datenqualität während der durchgeführten Evaluation. Zwei Aspekte fallen hierbei besonders schwer ins Gewicht. Zum einen weisen die aufgezeichneten Logs eine gewisse Anzahl an fehlerhaften oder vollständig leeren Gesprächen auf. Solche unbrauchbaren Einträge verzerren die statistische Verteilung der Scores und erschweren einen unverfälschten Vergleich der tatsächlichen Mittelwerte.

Zum anderen stellt die extrem niedrige Retrieval Abdeckung von lediglich 30\,\% ein massives Problem für die Beantwortung der zweiten Forschungsfrage dar. In mehr als zwei Dritteln aller simulierten Interaktionen konnte das RAG System keine passenden Kontextinformationen aus der bereitgestellten Dokumentendatenbank extrahieren. Das Sprachmodell musste folglich zwangsläufig auf sein internes parametrisches Wissen zurückgreifen, um die gestellten fachlichen Fragen zu beantworten. Da die RAG Pipeline in diesen Fällen faktisch inaktiv war, lief der Vergleich der verschiedenen Retrieval Konfigurationen weitgehend ins Leere. Die erfreulich hohe fachliche Korrektheit der Antworten ist somit primär der enormen Leistungsfähigkeit und dem immensen Vorwissen des zugrunde liegenden Sprachmodells zuzuschreiben und spiegelt nicht die Qualität des Information Retrievals wider.

\subsection{Limitationen (Stichprobengröße, Chunk Granularität, Judge Bias und Threshold Wahl)}

Die absolute Aussagekraft sowie die Übertragbarkeit der Studie werden durch mehrere methodische Entscheidungen und strukturelle Rahmenbedingungen limitiert. Zunächst ist die gewählte Stichprobengröße von lediglich 24 Szenarien pro Bedingung als verhältnismäßig gering einzustufen. Kleine Stichproben reduzieren die statistische Reliabilität erheblich und führen dazu, dass einzelne stark abweichende Dialoge die Gesamtergebnisse überproportional beeinflussen können.

Darüber hinaus liegt eine wesentliche technische Einschränkung in der verwendeten Chunk Granularität der Quelldokumente begründet. Die Art und Weise, wie die komplexen Laborunterlagen in durchsuchbare Textbausteine zerlegt wurden, passte offensichtlich nicht optimal zu den sehr spezifischen Anfragen der simulierten Nutzer. Diese unzureichende Segmentierung der zugrunde liegenden Daten ist die wahrscheinlichste Ursache für die zuvor beschriebene mangelhafte Retrieval Abdeckung.

Ebenso kritisch ist der sogenannte Judge Bias bei der Nutzung eines Sprachmodells als automatisierte Jury zu betrachten. Große Modelle neigen in ihrer Funktion als Evaluator systematisch dazu, besonders lange und stark formal gegliederte Antworten zu bevorzugen. Dieser Effekt erklärt die irreführend hohen Bewertungen der promptlosen Referenz im Bereich der schrittweisen Lernprogression, da diese Variante ungefragt umfassende Komplettlösungen in Listenform präsentierte.

Schließlich stellt die Festlegung des Schwellenwerts auf 0,7 eine rein konzeptionelle Threshold Wahl dar. Bereits eine geringfügige Anpassung dieser Bestehensgrenze würde die berechneten Erfolgsquoten signifikant verändern und könnte die daraus resultierende Rangfolge der evaluierten Promptvarianten unter Umständen entscheidend beeinflussen.